\chapter{Ziel der Arbeit}

Der Fokus der Arbeit gliedert sich in folgende Kernbereiche:

\begin{itemize}
    \item \textbf{Anforderungsanalyse und theoretische Modellierung:} Zu Beginn steht eine fundierte Grundlagenermittlung. Hierbei werden mithilfe von \textit{Motor-CAD} und ggf. \textit{ANSYS Maxwell} detaillierte theoretische Modelle erstellt, um die Verlustleistungen und thermischen Lastpfade der elektrischen Komponenten zu identifizieren.
    
    \item \textbf{Systemsimulation in Matlab/Simulink:} Um verschiedene Architekturen des Kühlkreislaufs (z. B. serielle vs. parallele Durchströmung) frühzeitig bewerten zu können, wird ein modulares Simulationsmodell in \textit{Matlab/Simulink} entwickelt. Dies ermöglicht eine quantitative Vergleichbarkeit unterschiedlicher Konzepte hinsichtlich ihrer Effizienz und thermischen Trägheit.
    
    \item \textbf{Detaillierte CFD-Untersuchungen:} Zur qualitativen Absicherung der Kühlgeometrien von Motor, Inverter und Akku werden numerische Strömungssimulationen (CFD) durchgeführt. Ziel ist es, lokale Hotspots zu identifizieren und die Wärmeabfuhr durch optimierte Kühlmantel-Designs zu maximieren.
    
    \item \textbf{Messtechnik und Validierung:} Ein wesentlicher Bestandteil der Arbeit ist die Auslegung der Sensorik (Durchfluss- und Temperaturmessung) sowohl für den fahrzeugseitigen Einsatz als auch für die Verifikation am Motorprüfstand. Die dort gewonnenen Messdaten dienen der Validierung der zuvor erstellten Simulationsmodelle.
    
    \item \textbf{Dokumentation und Parametervariation:} Abschließend erfolgt eine detaillierte Bewertung der Gesamtkühlarchitektur. Durch eine Sensitivitätsanalyse wird ermittelt, welche Parameter die Kühlleistung maßgeblich beeinflussen. Dies soll die zukünftige Auslegung vereinfachen und als Leitfaden für kommende Fahrzeuggenerationen dienen.
\end{itemize}

Durch diesen hybriden Ansatz aus Simulation und experimenteller Validierung wird sichergestellt, dass das Kühlsystem nicht nur theoretisch optimal dimensioniert, sondern auch unter realen Bedingungen belastbar ist.
